\documentclass[11pt]{article}
\usepackage[margin=1in]{geometry}
\usepackage{amsmath,amssymb,amsthm,mathtools}
\usepackage{booktabs}
\usepackage{microtype}
\usepackage[hidelinks]{hyperref}
\usepackage{enumitem}
\usepackage{array}

\newtheorem{theorem}{Theorem}
\newtheorem{lemma}[theorem]{Lemma}
\newtheorem{corollary}[theorem]{Corollary}
\newtheorem{proposition}[theorem]{Proposition}
\theoremstyle{definition}
\newtheorem{definition}[theorem]{Definition}
\newtheorem{remark}[theorem]{Remark}

\newcommand{\Z}{\mathbb{Z}}
\newcommand{\Fp}{\mathbb{F}_2}
\newcommand{\D}{\operatorname{D}}
\newcommand{\1}{\mathbf{1}}

\title{\textbf{Top-KPQ Completeness and Q-Space Geometry on Even Cyclic Rings}\\
Exact Collision Classes, Canonical Terminals, and Metric-Affine Symmetry}
\author{Muhammad Arshad}
\date{September 2026}

\begin{document}
\maketitle

\begin{abstract}
Let $R_H=\Z/(2H)\Z$ and let $c=(x_0,\ldots,x_{K-1})$ be a temporally ordered selection of $K$ ring positions. Its $m=K-1$ movements are $e_t=x_{t+1}-x_t$. We define the fingerprint
\[
F(c)=(K,P,Q),\qquad
P=x_{K-1}-x_0,\qquad
Q=\frac{L-\D(P)}2,
\]
where $L=\sum_t \D(e_t)$ and $\D$ is circular distance. Under the primitive same-sign repartitions
\[
\{+i,+j\}\leftrightarrow\{+(i-1),+(j+1)\},
\qquad
\{-i,-j\}\leftrightarrow\{-(i-1),-(j+1)\},
\]
together with temporal permutation, we prove that $(K,P,Q)$ is a complete invariant: two selections lie in the same collision class if and only if their fingerprints agree. Every class has a unique terminal described by four integers $(a,b,h,z)$. The same data induce an exact integer Q-space simplex, an exact minimum realizable order $K^\ast$, a $\Z_2$ quotient on the self-conjugate displacement sheets $P\in\{0,H\}$, and the metric-affine symmetry group $S_m\times V_4$. The MPRC specialization is $H=128$, i.e. $R_H=\Z_{256}$. Exhaustive computations are retained as regression evidence but are not substituted for the proofs.
\end{abstract}

\noindent\textbf{Keywords:} cyclic ring; discrete geometry; rewrite system; complete invariant; canonical form; Q-space; MPRC; $\Z_{256}$.

\medskip
\noindent\textbf{Reproducibility.}
The exact-integer reference gates are archived in the public \texttt{siliq-rotor} repository; no floating-point arithmetic is required for the theorem checks.

\section{Introduction}
This paper studies a finite collision geometry on the even cycle $R_H=\Z/(2H)\Z$. The motivating case is the MPRC ring $\Z_{256}$, but the principal statements are formulated for general $H$.

The symbol $K$ is not a tunable Top-K hyperparameter. It is the observed number of selected positions in a temporally ordered word. A $K$-position word contains
\[
m=K-1
\]
movements. The problem is to classify such movement data under a native local repartition law.

The central result is
\[
F(c_1)=F(c_2)\quad\Longleftrightarrow\quad c_1\sim c_2.
\]
Thus $(K,P,Q)$ does not merely survive the rewrites; it exactly labels their connected components.

\paragraph{Contributions.}
\begin{enumerate}[leftmargin=*]
\item We prove termination of the oriented primitive collision system by an integer potential.
\item We derive the unique terminal form $(a,b,h,z)$ and correct the slot count using $m=K-1$.
\item We prove that $(K,P,Q)$ is a complete invariant of the full rewrite congruence.
\item We derive the static Q-space simplex and its exact $m$-admissibility condition.
\item We derive the exact minimum selected-position order $K^\ast$.
\item We prove that the self-conjugate sheets $P=0,H$ quotient by the involution $\eta^+\leftrightarrow\eta^-$.
\item We classify the uniform metric-affine ring symmetries and obtain $S_m\times V_4$.
\end{enumerate}

\paragraph{Scope.}
Completeness is relative to the stated collision congruence. It does not identify byte equality, semantic equality, address equality, or any finer histogram metric with KPQ equality.

\section{Ring geometry and the frozen fingerprint}
For $e\in R_H$, define
\[
\D(e)=\min(e\bmod 2H,\,(-e)\bmod 2H),
\]
so $0\le \D(e)\le H$.

Let
\[
c=(x_0,\ldots,x_{K-1})
\]
be a temporally ordered selection and put
\[
m=K-1,\qquad e_t=x_{t+1}-x_t\pmod{2H}.
\]
Define
\[
P=\sum_{t=0}^{m-1}e_t=x_{K-1}-x_0\pmod{2H},
\]
\[
L=\sum_{t=0}^{m-1}\D(e_t),\qquad d=\D(P).
\]
On an even cycle, $L-d$ is even, so
\[
Q=\frac{L-d}{2}\in\Z_{\ge0}.
\]

\begin{definition}[Top-KPQ fingerprint]
The frozen fingerprint of $c$ is
\[
F(c)=(K,P,Q).
\]
Here $K$ counts selected positions, while $m=K-1$ counts movement slots.
\end{definition}

\section{Primitive collision relations}
Write ordinary movements as $+r$ or $-r$, $1\le r\le H-1$. The states $0$ and $H$ are sign-free. The primitive same-sign relations are
\[
\{+i,+j\}\leftrightarrow\{+(i-1),+(j+1)\},
\]
\[
\{-i,-j\}\leftrightarrow\{-(i-1),-(j+1)\},
\]
for $1\le i\le j\le H-1$ when the resulting magnitudes are legal. Temporal permutation of movement slots is also admitted.

Each primitive move preserves $m$, the signed sum $P$, and the length $L$. Hence it preserves $Q$.

\section{Termination}
Define
\[
\Psi(e_1,\ldots,e_m)=\sum_{r=1}^m \D(e_r)^2.
\]
Orient each primitive relation as
\[
(i,j)\longmapsto(i-1,j+1).
\]
Then
\[
(i-1)^2+(j+1)^2-i^2-j^2=2(j-i+1)\ge2.
\]

\begin{lemma}[Strict potential]
Every oriented primitive rewrite strictly increases $\Psi$.
\end{lemma}

Since $\D(e_r)\le H$,
\[
\Psi\le mH^2.
\]
Therefore:

\begin{corollary}[Termination]
Every movement multiset reaches an irreducible terminal after finitely many oriented primitive rewrites.
\end{corollary}

\section{The terminal shape}
An irreducible multiset cannot contain two ordinary positive magnitudes, since they would admit a positive repartition. The same holds on the negative side. Hence every terminal is of the form
\[
\{+a\}^{[a>0]}\cup\{-b\}^{[b>0]}\cup\{H\}^h\cup\{0\}^z,
\]
with $0\le a,b<H$.

Let $P$ be represented by its residue in $\{0,\ldots,2H-1\}$. Define
\[
a=\left(\frac{L+P}{2}\right)\bmod H,\qquad
b=\left(\frac{L-P}{2}\right)\bmod H.
\]
For a reachable class,
\[
h=\frac{L-a-b}{H}\in\Z_{\ge0}.
\]
The terminal counts movement slots, therefore
\[
\boxed{
z=m-h-[a>0]-[b>0]
=(K-1)-h-[a>0]-[b>0].
}
\]

\begin{remark}[Gate-discovered notation correction]
An earlier draft used $K$ rather than $m$ in the formula for $z$. That is off by one because $K$ counts positions and the terminal counts movements. The completeness theorem and the formulas for $(a,b,h)$ are unchanged.
\end{remark}

\begin{theorem}[Unique terminal]
For every reachable $(K,P,Q)$ class, the above $(a,b,h,z)$ is the unique irreducible terminal.
\end{theorem}

\begin{proof}
The invariants give $m=K-1$ and $L=2Q+\D(P)$. Any irreducible terminal has the stated shape. Its signed sum and total length force $a$ and $b$ modulo $H$; reachability forces the representatives in $[0,H)$ and makes $L-a-b$ a nonnegative multiple of $H$. This determines $h$, and the movement count then determines $z$.
\end{proof}

\section{Top-KPQ completeness}
\begin{theorem}[Top-KPQ Completeness]\label{thm:complete}
For two temporally ordered selections $c_1,c_2$,
\[
\boxed{
F(c_1)=F(c_2)\iff c_1\sim c_2.
}
\]
\end{theorem}

\begin{proof}
The reverse implication follows from invariance of $K,P,L$ under primitive rewrites and permutations.

Conversely, suppose $F(c_1)=F(c_2)$. Then both selections have the same $m=K-1$, the same $P$, and the same
\[
L=2Q+\D(P).
\]
By termination, both reduce to irreducible terminals. By the unique-terminal theorem, those terminals are identical. Since the generating collision relations are undirected in the congruence, both selections lie in the same component.
\end{proof}

Thus the canonicalization map
\[
\mathcal C(c)=T(F(c))
\]
is exact for the stated collision congruence.

\section{Static Q-space simplex}
Fix a reachable $(P,Q)$ class and write $d=\D(P)$. Choose the endpoint orientation. Let $A$ denote total ordinary magnitude in that direction, $B$ the total ordinary magnitude in the opposite direction, and $h$ the number of half-turn movements.

Define
\[
a=(Q+d)\bmod H,\qquad
r=\left\lfloor\frac{Q+d}{H}\right\rfloor,
\]
\[
b=Q\bmod H,\qquad
s=\left\lfloor\frac QH\right\rfloor,
\qquad
T=r+s.
\]
Every member of the class has
\[
A=a+H\eta^+,\qquad
B=b+H\eta^-,
\]
with
\[
\boxed{
\eta^++\eta^-+\eta_H=T,\qquad \eta_H=h.
}
\]
Hence the internal geometric coordinate lies on the integer simplex
\[
\Delta_T=\{(\eta^+,\eta^-,\eta_H)\in\Z_{\ge0}^3:
\eta^++\eta^-+\eta_H=T\}.
\]

For $S\ge0$, put
\[
c_H(S)=
\begin{cases}
0,&S=0,\\[3pt]
\left\lceil\dfrac{S}{H-1}\right\rceil,&S>0.
\end{cases}
\]

\begin{theorem}[Exact admissibility]
A simplex point is realizable with $m$ movement slots if and only if
\[
\boxed{
c_H(A)+c_H(B)+h\le m.
}
\]
\end{theorem}

\section{Minimum realizable order}
The canonical terminal uses
\[
m^\ast=T+[a>0]+[b>0]
\]
movement slots. Therefore the minimum selected-position count is
\[
\boxed{
K^\ast=1+T+[a>0]+[b>0].
}
\]

\section{Self-conjugate sheets}
Let
\[
C(e)=-e\pmod{2H}.
\]
Then $\D(Ce)=\D(e)$, so $L$ and $Q$ are fixed while $P\mapsto-P$. The fixed displacement sheets are exactly
\[
P\in\{0,H\}.
\]
On either sheet, the base residues satisfy $a=b$. Conjugation swaps the oriented totals $A$ and $B$ and leaves the half-turn count fixed. Therefore
\[
(\eta^+,\eta^-,\eta_H)\mapsto(\eta^-,\eta^+,\eta_H).
\]

\begin{theorem}[Unoriented sheet quotient]
On $P=0$ and $P=H$, the unoriented static Q-space is exactly
\[
\boxed{\Delta_T/\Z_2},
\]
where the nontrivial element swaps $\eta^+$ and $\eta^-$.
\end{theorem}

For the unrestricted simplex, Burnside gives
\[
\left|\Delta_T/\Z_2\right|
=
\frac12\left[
\binom{T+2}{2}
+\left\lfloor\frac T2\right\rfloor+1
\right].
\]

\section{Metric-affine symmetry}
Consider a uniform affine ring automorphism
\[
f(e)=ae+b\pmod{2H},\qquad \gcd(a,2H)=1.
\]
Call $f$ metric-affine if globally either
\[
\D(f(e))=\D(e)
\]
or
\[
\D(f(e))=H-\D(e).
\]

\begin{theorem}[Exact ring symmetry]
The metric-affine ring symmetry group is exactly
\[
\boxed{
V_4=\{I,C,A,AC\},
}
\]
where
\[
C(e)=-e,\qquad A(e)=e+H.
\]
\end{theorem}

\begin{proof}
In the preserving case, $e=0$ forces $b=0$ and $e=1$ forces $\D(a)=1$, hence $a=\pm1$. In the dual case, $e=0$ forces $b=H$; using $\D(x+H)=H-\D(x)$ at $e=1$ again gives $\D(a)=1$, hence $a=\pm1$.
\end{proof}

The uniform $V_4$ action commutes with arbitrary permutations of the $m$ movement coordinates. Therefore:

\begin{corollary}[Full stated metric-affine group]
\[
\boxed{
G_m=S_m\times V_4.
}
\]
\end{corollary}

Define
\[
\Lambda=2L-mH.
\]
Then
\[
C:\quad P\mapsto-P,\quad L\mapsto L,\quad \Lambda\mapsto\Lambda,
\]
while
\[
A:\quad P\mapsto P+mH,\quad L\mapsto mH-L,\quad \Lambda\mapsto-\Lambda.
\]

\section{Verification}
The proof is analytical; finite computation serves as falsification and regression. The reference programs exhaust:
\begin{itemize}[leftmargin=*]
\item rewrite components versus $(K,P,Q)$ classes on small rings;
\item strict potential increase;
\item unique closed-form terminals;
\item exact simplex coverage and $K^\ast$;
\item the self-conjugate $\Z_2$ action;
\item every affine automorphism in the tested rings.
\end{itemize}

For the MPRC ring $H=128$, every two-movement multiset is checked:
\[
\binom{257}{2}=32{,}896.
\]
Under the corrected convention this is the $K=3$ selected-position case.

\section{Computational interpretation}
Theorem~\ref{thm:complete} permits exact collision-class canonicalization. It does not erase finer equality relations. A runtime may use KPQ as an exact narrowing key only when the indexed equivalence is the collision congruence; exact byte or record equality must still decide individual candidates when required. Likewise, a metric defined only on a finer histogram fiber must not be extended across an entire KPQ class.

\section{Limitations and open questions}
This paper does not claim that KPQ equality is semantic equality, address equality, or universal identity. It does not identify the $252\to576$ arithmetic transition with a CXR grid lift. It also does not address the dyadic ramification of the collision algebra; that is treated separately.

\section{Conclusion}
The native even-cycle collision system admits an exact finite classification. The fingerprint
\[
(K,P,Q)
\]
is complete, every class has a unique terminal, the internal multiplicity has an exact simplex geometry, the self-conjugate sheets quotient by a single involution, and the metric-affine symmetry is exactly
\[
S_m\times V_4.
\]
No ranking score, fitted threshold, or probability simplex is required for this classification.

\section*{Reproducibility record}
The frozen theorem register is:
\begin{center}
\url{https://github.com/muhammadarshad/pages-mprc/tree/master/research/top-kpq}
\end{center}
The collision and symmetry gates are archived in:
\begin{center}
\url{https://github.com/muhammadarshad/siliq-rotor/tree/master/model/motif}
\end{center}

\begin{thebibliography}{9}
\bibitem{ArshadSimplex}
M. Arshad,
\emph{Quantinion Simplex Cubic Decomposition over 2-Power Rings: Global Identifiability from Binary Contractions and Exact Bit Lifting},
Zenodo, 2026.
doi:10.5281/zenodo.22850767.

\bibitem{ArshadCapacity}
M. Arshad,
\emph{Exact Free Capacity of Symmetric Rank-One Cubic Dictionaries over Dyadic Rings},
Zenodo, 2026.
doi:10.5281/zenodo.22850810.

\bibitem{ArshadOperator}
M. Arshad,
\emph{An Eight-State Operator and Transpose Geometry over Dyadic Rings: Trace-Zero Modules, Mixed-Radix Addressing, and Tensor-Lift Kernels},
Zenodo, 2026.
doi:10.5281/zenodo.22850830.
\end{thebibliography}

\end{document}
